\begin{description}
\item[(a)] Concentrations of electrons, $N_e$, and He nuclei, $N_{He}$, is
  related to the concentration of H nuclei, $N_H$:
  \begin{equation*}
  N_{He} = 0.1\,N_H
  \tag*{(1 point)}
  \end{equation*}
  \begin{equation*}
  N_e = N_H + 2N_{He} = 1.2\,N_H
  \tag*{(2 points)}
  \end{equation*}
  The total X-ray emission is a sum of the bremmstrahlung generated by
  % sic: "bremmstrahlung" as in the original
  interaction of electrons with H and He nuclei:
  \[
  L = 6 \cdot 10^{-41}\, N_e\, T^{\frac{1}{2}} V \left(N_H + Z_{He}^2 N_{He}\right),
  \]
  where $Z_{He} = 2$. The luminosity $L$ is expressed by the concentration of
  H nuclei:
  \begin{gather*}
  L = 6 \cdot 10^{-41}\; 1.2\,N_H\; T^{\frac{1}{2}}\, V\; 1.4\,N_H, \\
  \Longrightarrow\;
  L = (10.08 \times 10^{-41})\, T^{\frac{1}{2}} V N_H^2
  \approx 10^{-40}\, T^{\frac{1}{2}} V N_H^2.
  \tag*{(3 points)}
  \end{gather*}
  The volume $V$ is given by:
  \begin{equation*}
  V = \frac{4}{3}\pi R^3 = 1.54 \cdot 10^{67}~\mathrm{m^3}.
  \tag*{(3 points)}
  \end{equation*}
  Thus, the concentration of $N_H$:
  \begin{equation*}
  N_H = \left(\frac{L}{10^{-40}\, T^{\frac{1}{2}} V}\right)^{1/2}
  \approx 8.48 \times 10^{2}~\mathrm{m^{-3}}.
  \tag*{(3 points)}
  \end{equation*}
  To obtain the total mass of the plasma, $M$, one should multiply the volume
  $V$ by the the sum of H and He mass densities:
  % sic: "the the" as in the original
  \begin{equation*}
  M = V\left(N_H m_H + N_{He} m_{He}\right)
  = 3.03 \times 10^{43}~\mathrm{kg}
  \approx 1.52 \times 10^{13}\,M_\odot,
  \tag*{(4 points)}
  \end{equation*}
  where $m_H$ and $m_{He}$ are the masses of hydrogen and helium atoms.

\item[(b)] Let $L$ be the mean free path of a photon. The number of electrons
  in a cylinder with a cross section area of $\sigma$ and a length of $L$
  equals
  \[
  N = nL\sigma = 1 \;\Longrightarrow\; L = 1/(n\sigma).
  \]
  Assuming $n = N_e = 1.2 N_H = 1.018 \times 10^{3}~\mathrm{m^{-3}}$, we get
  \begin{equation*}
  L = \frac{1}{1.018 \times 10^{3} \cdot 6.65 \times 10^{-29}}
  = 1.48 \times 10^{25}~\mathrm{m}
  = 4.79 \times 10^{8}~\mathrm{pc}
  \approx 500~\mathrm{Mpc}.
  \tag*{(5 points)}
  \end{equation*}
  (The radius of the cluster is $\sim 1000$ times smaller than $L$, so the
  interactions between the CMB photons and hot electrons are rare.)

\item[(c)] Using Wien's law,
  \begin{gather*}
  \lambda = \frac{b}{T} = \frac{2.898 \times 10^{-3}}{2.73}
  = 1.07 \times 10^{-3}~\mathrm{m} \\
  \Longrightarrow\; E_0 = \frac{hc}{\lambda} = 1.85 \times 10^{-22}~\mathrm{J}
  \tag*{(4 points)}
  \end{gather*}
  % sic: parts (c)+(d) point values in the source solution (4+3+2) do not
  % match the statement's split (3+6); transcribed as printed.

\item[(d)] We first have to estimate the typical velocity $v_e$ of electrons
  using the formula for the kinetic energy of the particles in a gas.
  \[
  \frac{1}{2}mv^2 = \frac{3}{2}kT
  \;\Longrightarrow\;
  v_e = \sqrt{\frac{3kT}{m_e}}
  \]
  \begin{equation*}
  \therefore\; v_e
  = \sqrt{\frac{3 \cdot 1.381 \times 10^{-23} \cdot 8 \times 10^{7}}
               {9.109 \times 10^{-31}}}
  = 6.032 \times 10^{7}~\mathrm{m\,s^{-1}} = 0.202c
  \tag*{(3 points)}
  \end{equation*}
  The energy of upscattered photons is:
  \begin{equation*}
  E' = \frac{1+\beta}{1-\beta}\,E_0 \approx 1.5 E_0
  = 2.78 \times 10^{-22}~\mathrm{J}.
  \tag*{(2 points)}
  \end{equation*}
  \emph{Note:} The formula for the kinetic energy of particles in an ideal gas
  remains valid even at such high temperatures. This is because
  \[
  x = \frac{m_e c^2}{kT}
  = \frac{9.109 \times 10^{-31} \cdot (2.998 \times 10^{8})^2}
         {1.381 \times 10^{-23} \cdot 8 \times 10^{7}}
  = 74.1 \gg 1,
  \]
  and the electrons can be treated as non-relativistic. The full relativistic
  formula for the rms velocity of particles in an ideal gas is
  \[
  \langle v_e^2 \rangle = \frac{xc^2}{K_2(x)}
  \int_0^\infty \frac{\sinh^4\varphi}{\cosh\varphi}\,
  e^{-x\cosh\varphi}\,d\varphi,
  \]
  where $K_2(x)$ is a modified Bessel function of the second kind. It can be
  shown that for $x = 74.1$ this formula yields
  $\sqrt{\langle v_e^2 \rangle} = 0.198c$, which is virtually identical to the
  prediction of the non-relativistic formula.
\end{description}
